\textbf{Executive overview.} Vision-Language-Action (VLA) models are large neural policies that map (RGB image, language instruction) to robot actions in a single forward pass through a vision-language backbone. The class crystallized in July 2023 with RT-2 {[}Brohan et al., 2023{]}, scaled rapidly through OpenVLA-7B (June 2024, 76.5\% average on BridgeData/RT-1), Octo (RSS 2024, 27M/93M diffusion), π0 (October 2024, flow-matching at 50 Hz), CogACT (November 2024, +55\% real-world over OpenVLA), SpatialVLA, X-VLA, Gemini Robotics 1.5, Helix, π0.5, Xiaomi-Robotics-0, and Toyota's Large Behavior Models in \emph{Science Robotics} {[}Barreiros et al., 2026{]}. The taxonomy stabilizes along four axes --- VLM backbone (PaLI, Llama, Qwen2-VL, PaliGemma), action head (256-bin discrete, MLP, diffusion, flow matching), system topology (monolithic, dual-system, modular), reasoning style (end-to-end, ECoT, Code-as-Policy, affordance) --- and most state-of-the-art systems by 2026 combine a 3--7B backbone with a continuous head and a slow-fast topology. Pretraining is dominated by Open X-Embodiment (22 embodiments, 527 skills, \textasciitilde1M trajectories, RLDS format), augmented by DROID (76k trajectories, 564 scenes, 86 tasks), RoboMIND (107k, 4 embodiments, 479 tasks), Mobile ALOHA, and retargeted human video. Evaluation has migrated from per-paper demos to LIBERO/CALVIN/SimplerEnv simulators and to distributed real-world services AutoEval {[}Zhou et al., 2025{]} and RoboArena {[}Atreya et al., 2025{]}, with mid-2026 reference \ldots{}
